\documentclass[../../main.tex]{subfiles}

\begin{document}

\chapter[Analisis Leksikal]{Analisis Leksikal: Token, Lexeme, Pattern,\\ dan Ekspresi Reguler}
\label{chap:leksikal}

\begin{outcomebox}
\textbf{Sub-CPMK:} \subcpmk{2.1}\\[3pt]
Setelah menyelesaikan bab ini, mahasiswa mampu:
\begin{enumerate}[label=\arabic*.,leftmargin=*]
  \item menjelaskan peran dan antarmuka lexer dalam pipeline
        kompilasi;
  \item membedakan token, lexeme, dan pattern beserta atributnya;
  \item mendefinisikan dan mengoperasikan ekspresi reguler;
  \item menulis spesifikasi token dengan prioritas dan aturan
        \eng{maximal munch};
  \item menyusun spesifikasi token Subset-C dan header \file{token.h}.
\end{enumerate}
\end{outcomebox}

\section*{Materi Pokok}
\begin{itemize}[leftmargin=*]
  \item Peran lexer dan antarmukanya dengan parser.
  \item Token, lexeme, pattern, dan atribut token.
  \item Ekspresi reguler: definisi, operasi, dan identitas.
  \item Notasi, prioritas aturan, dan maximal munch.
  \item Spesifikasi token Subset-C dan penanganan kesalahan leksikal.
\end{itemize}

\subfile{section-01}
\subfile{section-02}
\subfile{section-03}
\subfile{section-04}
\subfile{section-05}

\section*{Contoh}
\addcontentsline{toc}{section}{Contoh}

\begin{contoh}[Menokenisasi secara manual]
Diberikan program
\begin{lstlisting}[language=C]
int x = 10;
x = x + 1;
\end{lstlisting}
Token yang dihasilkan:
\begin{center}
\small
\begin{tabular}{@{}lll@{}}
\toprule
\textbf{Baris:Kolom} & \textbf{Token} & \textbf{Lexeme} \\
\midrule
1:1 & \code{TOK\_INT} & \code{int} \\
1:5 & \code{TOK\_IDENT} & \code{x} \\
1:7 & \code{TOK\_ASSIGN} & \code{=} \\
1:9 & \code{TOK\_INT\_LIT} & \code{10} \\
1:11 & \code{TOK\_SEMI} & \code{;} \\
2:1 & \code{TOK\_IDENT} & \code{x} \\
2:3 & \code{TOK\_ASSIGN} & \code{=} \\
2:5 & \code{TOK\_IDENT} & \code{x} \\
2:7 & \code{TOK\_PLUS} & \code{+} \\
2:9 & \code{TOK\_INT\_LIT} & \code{1} \\
2:10 & \code{TOK\_SEMI} & \code{;} \\
3:1 & \code{TOK\_EOF} & --- \\
\bottomrule
\end{tabular}
\end{center}
Perhatikan bahwa spasi tidak menghasilkan token, dan setiap token
memiliki lokasi awal.
\end{contoh}

\begin{contoh}[Token, lexeme, dan pattern]
Untuk token \code{FLOAT\_LIT}, lexeme \code{3.14} cocok dengan
pattern \verb![0-9]+\.[0-9]+!. Lexeme \code{42} bukan anggota token
ini karena tidak memiliki titik. Sebaliknya, \code{42} adalah lexeme
untuk token \code{INT\_LIT}. Pattern mendefinisikan himpunan lexeme;
token adalah nama himpunan tersebut.
\end{contoh}

\begin{contoh}[Kata kunci versus identifier]
Misalkan tabel kata kunci berisi \code{int}, \code{if}, dan
\code{while}. Masukan \code{int x;} menghasilkan token
\code{INT}, \code{IDENT(x)}, \code{SEMI}. Masukan \code{integer x;}
menghasilkan \code{IDENT(integer)}, \code{IDENT(x)}, \code{SEMI},
karena \code{integer} tidak ada dalam tabel kata kunci. Lexer tidak
menebak berdasarkan awalan; ia mencocokkan seluruh lexeme.
\end{contoh}

\begin{contoh}[Maximal munch]
Pada masukan \code{a<=b}, lexer membaca \code{<} lalu melihat
karakter berikutnya \code{=}. Karena \code{<=} cocok dan lebih
panjang daripada \code{<}, lexer memilih \code{LE}. Pada masukan
\code{a<<b}, tidak ada operator \code{<<} dalam Subset-C, sehingga
lexer memilih \code{LT} dua kali. Aturan maximal munch bekerja pada
pattern yang tersedia, bukan pada operator yang mungkin ada dalam
bahasa lain.
\end{contoh}

\begin{contoh}[Kesalahan leksikal]
Masukan \code{int a = 3.;} mengandung float tanpa digit setelah
titik. Lexer menghasilkan \code{TOK\_ERROR} pada posisi titik,
mencatat pesan \code{baris 1, kolom 10: digit diharapkan setelah
titik}, lalu melanjutkan. Parser menerima token kesalahan dan dapat
melaporkan masalah tanpa menghentikan seluruh kompilasi.
\end{contoh}

\section*{Aktivitas Pembelajaran dan Praktikum}
\addcontentsline{toc}{section}{Aktivitas Pembelajaran dan Praktikum}

\begin{table}[H]
\centering
\caption{Rencana aktivitas Bab V}
\label{tab:ch05-aktivitas}
\begin{tabularx}{\textwidth}{@{}llX@{}}
\toprule
\textbf{Sesi} & \textbf{Bentuk} & \textbf{Kegiatan} \\
\midrule
1 & Ceramah & Peran lexer dan antarmuka token stream \\
2 & Latihan & Menokenisasi program secara manual \\
3 & Diskusi & Ekspresi reguler dan prioritas aturan \\
4 & Praktikum & Menyusun spesifikasi token dan kasus uji \\
\bottomrule
\end{tabularx}
\end{table}

\subsection*{P1-A: Spesifikasi Token dan 30 Kasus Uji}
\textbf{Tujuan.} Menyusun spesifikasi token Subset-C yang lengkap
sebelum implementasi lexer.

\textbf{Langkah kerja.}
\begin{enumerate}[label=\arabic*.,leftmargin=*]
  \item Lengkapi \file{spec/tokens.md} dengan tabel pattern untuk
        setiap token.
  \item Tuliskan daftar kata kunci dan aturan escape untuk char serta
        string.
  \item Buat 30 kasus uji di \file{tests/lexer/}: 10 kasus normal,
        10 kasus batas, dan 10 kasus kesalahan.
  \item Untuk setiap kasus, tuliskan berkas \code{.sc} dan berkas
        \code{.expected} dengan format \code{BARIS:KOLOM JENIS LEXEME}.
  \item Periksa bahwa urutan aturan memenuhi maximal munch.
\end{enumerate}

\textbf{Definisi selesai.}
\begin{checklist}
  \item Semua token memiliki pattern.
  \item 30 kasus uji tersedia dan terdokumentasi.
  \item Format keluaran tes konsisten.
  \item Tidak ada aturan yang saling bertentangan.
\end{checklist}

P1-A belum dinilai secara sumatif. Hasilnya menjadi dasar Praktikum 1
(8\%) pada Bab VII.

\section*{Latihan dan Refleksi}
\addcontentsline{toc}{section}{Latihan dan Refleksi}

\subsection*{Latihan}
\begin{enumerate}[label=\arabic*.,leftmargin=*]
  \item Jelaskan perbedaan token, lexeme, dan pattern dengan contoh.
  \item Tuliskan pattern untuk identifier yang boleh mengandung garis
        bawah dan angka, tetapi tidak boleh diawali angka.
  \item Tuliskan pattern untuk bilangan pecahan tanpa bagian
        eksponen.
  \item Mengapa operator dua karakter harus ditempatkan sebelum
        operator satu karakter?
  \item Berikan contoh di mana maximal munch mencegah kesalahan
        parsing.
  \item Rancang aturan untuk komentar blok yang tidak bersarang.
  \item Jelaskan perilaku lexer ketika menemukan string yang tidak
        ditutup.
  \item Bandingkan strategi kata kunci sebagai identifier yang
        diperiksa tabel dengan strategi kata kunci sebagai pattern
        tersendiri.
\end{enumerate}

\subsection*{Refleksi}
\begin{reflectionbox}
\begin{enumerate}[label=\arabic*.,leftmargin=*]
  \item Bagian mana dari spesifikasi token yang paling sulit Anda
        tulis? Mengapa?
  \item Apakah Anda dapat menjelaskan mengapa lexer tidak perlu
        mengetahui tipe variabel?
  \item Kasus batas apa yang sering Anda lupakan saat menulis
        spesifikasi?
\end{enumerate}
\end{reflectionbox}

\section*{Asesmen}
\addcontentsline{toc}{section}{Asesmen}

Asesmen Bab V menjadi bagian Tugas Teori (8\%) dan prasyarat
Praktikum 1 (8\%). Fokusnya adalah ketepatan spesifikasi token dan
pemahaman maximal munch.

\begin{table}[H]
\centering
\caption{Instrumen asesmen Bab V}
\label{tab:ch05-asesmen}
\begin{tabularx}{\textwidth}{@{}lXl@{}}
\toprule
\textbf{Instrumen} & \textbf{Indikator} & \textbf{Bobot} \\
\midrule
Tugas Teori & Pattern token, prioritas, dan contoh benar & 8\% \\
P1-A & Spesifikasi token dan 30 kasus uji lengkap & prasyarat \\
Kuis & Token, lexeme, pattern, dan maximal munch & formatif \\
\bottomrule
\end{tabularx}
\end{table}

Kriteria kelulusan: mahasiswa dapat menulis spesifikasi token yang
tidak ambigu, menjelaskan prioritas aturan, dan menghasilkan kasus uji
yang mencakup kasus normal, batas, dan kesalahan.

\section*{Checklist Pencapaian Kompetensi}
\addcontentsline{toc}{section}{Checklist Pencapaian Kompetensi}

\begin{checklist}
  \item Saya dapat menjelaskan peran lexer dalam pipeline.
  \item Saya dapat membedakan token, lexeme, dan pattern.
  \item Saya dapat menyebutkan atribut yang dibawa token.
  \item Saya dapat membuat tabel ekspresi reguler untuk setiap
        kategori token Subset-C.
  \item Saya dapat menjelaskan prioritas operator ekspresi reguler.
  \item Saya dapat menerapkan aturan maximal munch.
  \item Saya dapat menyusun aturan prioritas token.
  \item Saya dapat menangani kesalahan leksikal tanpa menghentikan
        proses.
  \item Saya sudah menyelesaikan P1-A: tabel ekspresi reguler dan 30
        kasus uji.
\end{checklist}

\section*{Rangkuman}
\addcontentsline{toc}{section}{Rangkuman}

Pada bab ini telah dibahas:
\begin{itemize}[leftmargin=*]
  \item peran lexer dan antarmuka pull model dengan parser;
  \item perbedaan token, lexeme, pattern, dan atribut token;
  \item definisi, operasi, prioritas, dan identitas ekspresi reguler;
  \item aturan maximal munch dan prioritas spesifikasi token;
  \item tabel ekspresi reguler untuk setiap kategori token;
  \item spesifikasi token Subset-C, header \file{token.h}, dan
        penanganan kesalahan leksikal.
\end{itemize}

Bab VI akan menghubungkan ekspresi reguler dengan mesin pengenal:
NFA, DFA, konversi NFA ke DFA, dan minimalisasi DFA. Teori ini menjadi
dasar implementasi lexer pada Bab VII.

\section*{Evaluasi}
\addcontentsline{toc}{section}{Evaluasi}

\begin{enumerate}[label=\arabic*.,leftmargin=*]
  \item Definisikan token, lexeme, dan pattern. Berikan satu contoh
        untuk masing-masing.
  \item Tuliskan ekspresi reguler untuk identifier Subset-C.
  \item Tuliskan ekspresi reguler untuk bilangan pecahan sederhana.
  \item Jelaskan aturan maximal munch dan berikan dua contoh.
  \item Mengapa urutan aturan dalam spesifikasi token penting?
  \item Rancang aturan untuk komentar baris dan komentar blok.
  \item Diberikan masukan \code{a<=b==c}, tuliskan aliran token yang
        dihasilkan.
  \item Jelaskan bagaimana lexer harus menangani karakter ilegal.
  \item Buat tabel ekspresi reguler untuk setiap kategori token
        Subset-C dan berikan satu contoh lexeme.
\end{enumerate}

\subsection*{Kunci dan Rubrik Singkat}
\begin{itemize}[leftmargin=*]
  \item Soal 1--3: ketepatan definisi dan pattern (70\%),
        contoh (30\%).
  \item Soal 4--6: pemahaman aturan dan prioritas (70\%),
        kelengkapan (30\%).
  \item Soal 7--9: kebenaran token, tabel, dan perilaku (80\%),
        penjelasan (20\%).
\end{itemize}


\ifSubfilesClassLoaded{%
  \bibliographystyle{apalike}%
  \bibliography{../../references}%
}{}

\end{document}
